\documentclass[notes=show]{beamer}
%%%%%%%%%%%%%%%%%%%%%%%%%%%%%%%%%%%%%%%%%%%%%%%%%%%%%%%%%%%%%%%%%%%%%%%%%%%%%%%%%%%%%%%%%%%%%%%%%%%%%%%%%%%%%%%%%%%%%%%%%%%%%%%%%%%%%%%%%%%%%%%%%%%%%%%%%%%%%%%%%%%%%%%%%%%%%%%%%%%%%%%%%%%%%%%%%%%%%%%%%%%%%%%%%%%%%%%%%%%%%%%%%%%%%%%%%%%%%%%%%%%%%%%%%%%%
\AtBeginSubsection[
  {\frame<beamer>{\frametitle{Outline}   
    \tableofcontents[currentsection,currentsubsection]}}%
]%
{
  \frame<beamer>{ 
    \frametitle{Outline}   
    \tableofcontents[currentsection,currentsubsection]}
}
\usepackage{ulem}
\usepackage{pdflscape}
\setbeamertemplate{caption}{\insertcaption}
\usepackage{color,xcolor,ucs}
\usepackage{beamerthemesplit}
\usepackage{graphicx}
\usepackage{amsmath}
\usepackage{adjustbox}
\usepackage{pdfpages} 
\definecolor{mintgreen}{rgb}{0.6, 1.0, 0.6}
\makeatletter
\@addtoreset{subfigure}{framenumber}% subfigure counter resets every frame
\makeatother
\newcommand*{\Scale}[2][4]{\scalebox{#1}{$#2$}}%
\newcommand*{\Resize}[2]{\resizebox{#1}{!}{$#2$}}%
\usepackage{amsfonts}
\usepackage{booktabs}
\usepackage{color,soul}
\usepackage{pdfpages}
\setboolean{@twoside}{false}
\usepackage{amsmath}
\usepackage{pdflscape} 
\usepackage{eurosym}
\usepackage{amstext} % for \text
\DeclareRobustCommand{\officialeuro}{%
  \ifmmode\expandafter\text\fi
  {\fontencoding{U}\fontfamily{eurosym}\selectfont e}}
\usepackage[caption = false]{subfig}
\usepackage{appendixnumberbeamer} 
\usepackage{graphicx}
\usepackage{mathpazo}
\usepackage{hyperref}
\usepackage{multimedia}
\usepackage{graphicx}
\usepackage{multirow}
\usepackage{subfig}
\usepackage{verbatim}
\usepackage{booktabs, multicol, multirow}
\setcounter{MaxMatrixCols}{10}
\input{Macros.tex}
\AtBeginSection[]
{
  \begin{frame}[noframenumbering]
    \frametitle{Outline for section \thesection}
    \tableofcontents[currentsection]
  \end{frame}
}


\usepackage{xcolor,soul}
\renewcommand<>{\hl}[1]{\only#2{\beameroriginal{\hl}}{#1}}

% https://tex.stackexchange.com/questions/41683/why-is-it-that-coloring-in-soul-in-beamer-is-not-visible
\makeatletter
\newcommand\SoulColor{%
  \let\set@color\beamerorig@set@color
  \let\reset@color\beamerorig@reset@color}
\makeatother
\SoulColor
\AtBeginSection{\frame{\sectionpage}}

\newsavebox\mybox
\newenvironment{aquote}[1]
  {\savebox\mybox{#1}\begin{quote}\openautoquote\hspace*{-.7ex}}
  {\unskip\closeautoquote\vspace*{1mm}\signed{\usebox\mybox}\end{quote}}

\newtheorem{proposition}[theorem]{Proposition}
\newtheorem{Assumption}[theorem]{Assumption}
\newenvironment{stepenumerate}{\begin{enumerate}[<+->]}{\end{enumerate}}
\newenvironment{stepitemize}{\begin{itemize}[<+->]}{\end{itemize} }
\newenvironment{stepenumeratewithalert}{\begin{enumerate}[<+-| alert@+>]}{\end{enumerate}}
\newenvironment{stepitemizewithalert}{\begin{itemize}[<+-| alert@+>]}{\end{itemize} }
\usetheme{Madrid}
\makeatletter
\setbeamertemplate{footline}
{
  \leavevmode%
  \hbox{%
  \begin{beamercolorbox}[wd=.333333\paperwidth,ht=2.25ex,dp=1ex,center]{author in head/foot}%
    \usebeamerfont{author in head/foot}\insertshortauthor~~\beamer@ifempty{\insertshortinstitute}{}{(\insertshortinstitute)}
  \end{beamercolorbox}%
  \begin{beamercolorbox}[wd=.333333\paperwidth,ht=2.25ex,dp=1ex,center]{title in head/foot}%
    \usebeamerfont{title in head/foot}\insertshorttitle
  \end{beamercolorbox}%
  \begin{beamercolorbox}[wd=.333333\paperwidth,ht=2.25ex,dp=1ex,right]{date in head/foot}%
    \usebeamerfont{date in head/foot}\insertshortdate{}\hspace*{2em}
%    \insertframenumber{} / \inserttotalframenumber\hspace*{2ex} % DELETED
  \end{beamercolorbox}}%
  \vskip0pt%
}
\usepackage[style=british]{csquotes}

\def\signed #1{{\leavevmode\unskip\nobreak\hfil\penalty50\hskip1em
  \hbox{}\nobreak\hfill #1%
  \parfillskip=0pt \finalhyphendemerits=0 \endgraf}}

\makeatother
\usepackage{pdfpages}
\begin{document}
	\title[]{\large Wage Inequality Wars - Part II 		\vspace{0.5cm}
	}
	\author[Raffaele Saggio - Econ 560]{Raffaele Saggio}
	\institute[]{UBC}
	\date{\today \\}
	\maketitle

%%%%%%%%%%%%%%%%%%%%%%%%%	
%%%%%%%%%%%%%%%%%%%%%%%%%
%%%%%%%%%%%%%%%%%%%%%%%%%
%%%%%%%%%%%%%%%%%%%%%%%%%
\setbeamertemplate{footline}[frame number]{}

%gets rid of bottom navigation symbols
\setbeamertemplate{navigation symbols}{}

%gets rid of footer
%will override 'frame number' instruction above
%comment out to revert to previous/default definitions
\setbeamertemplate{footline}{}
%%%
\begin{frame}
\frametitle{Today}
\begin{itemize}
\item End of 1990's: consensus on importance of skill biased technological change (SBTC).
\medskip
 \begin{aquote}{Card and Di Nardo, 2002}
      The recent inequality literature reaches virtually unanimous agreement that the relative demand for highly skilled workers increased in the 1980s causing earnings inequality to increase.
      \end{aquote}
     \pause 
\bigskip
\item Mid 2000's: open war! East-West coast divide 
\bigskip
\item Key arguments for the critique: institutions, within group inequality, gender wage inequality, racial wage inequality.
\end{itemize}
\end{frame}
%%%
\begin{frame}
\frametitle{East vs. West Coast Battles, Labor Economics style}
\begin{figure}[htb]
\centering
\includegraphics[width=1\textwidth]{aux/east_coast_}
\end{figure}
\end{frame}
\begin{frame}
\frametitle{Card and Di Nardo 2002}
\begin{itemize}
\item Original title ``10 Things we hate about SBTC"
\medskip
\item SBTC provides an incomplete and sometimes undisciplined (``name the residual") account of changes in wage structure.
\medskip
\item Highlight: \underline{Problems} (patterns in the data inconsistent with SBTC) \underline{Puzzles} (patterns that could be explained by SBTC but seemed to be due to other factors)
\medskip
\pause
\item Main message: economists love uni-causal stories but in real world is really hard to get ``uni-causal" implications.
\end{itemize}
\end{frame}
%%
\begin{frame}
\frametitle{More moments: 1990s era, flat period of inequality?}
\begin{figure}[htb]
\centering
\includegraphics[width=1\textwidth]{aux/cdn00}
\end{figure}
\end{frame}
\begin{frame}
\frametitle{Were the 1980's an aberration?}
\begin{figure}[htb]
\centering
\includegraphics[width=1\textwidth]{aux/cdnBASE}
\end{figure}
\end{frame}
%%
\begin{frame}
\frametitle{Technology or Tautology?}
\begin{figure}[htb]
\centering
\includegraphics[width=1\textwidth]{aux/cdn1}
\end{figure}
\end{frame}
%%%
\begin{frame}
\frametitle{Several puzzles}
\begin{figure}[htb]
\centering
\includegraphics[width=0.82\textwidth]{aux/cdn4}
\end{figure}
\end{frame}
%%
\begin{frame}
\frametitle{Conclusions}
\begin{itemize}
\item 76\% of the rise in the 90-10 wage gap b/w 1979-1999 occurred between 1979-1984 $\rightarrow$ first year that the IBM-AT was introduced.
\medskip
\item SBTC may have played a role but it is unlikely to be the only story (too many problems/puzzles associated with that).
\medskip
\pause
\item If it's not SBTC then what is it? 
\end{itemize}
\end{frame}
%%
%%
\begin{frame}
\frametitle{Quite a coincidence!}
\begin{figure}[htb]
\centering
\includegraphics[width=1\textwidth]{aux/cdn5}
\end{figure}
\end{frame}
%%
%%
\begin{frame}
\frametitle{Quite a coincidence!}
\begin{figure}[htb]
\centering
\includegraphics[width=1\textwidth]{aux/cdn6}
\end{figure}
\end{frame}
%%
%%
\begin{frame}
\frametitle{Lee (1999)}
\begin{figure}[htb]
\centering
\includegraphics[width=0.9\textwidth]{aux/lee0}
\end{figure}
\end{frame}
%%
%%
%%
\begin{frame}
\frametitle{Card and Lemieux 2001}
\begin{itemize}
\item Original Title: ``It's supply, stupid". 
\bigskip
\pause
\item  Starting point: Are we sure we really need technical change to explain the 80s rise in return to schooling?
 \bigskip
 \item What if we are just observing a reduction in relative supply of college grads among recent cohorts?
  \bigskip
  \item In order for the last point to have a ``bite" we need to have imperfect substitution across age groups within education category.
   \bigskip
   \item International comparison (USA + UK + Canada).
\end{itemize}
\end{frame}
%%
\begin{frame}
\frametitle{Starting observation}
\begin{figure}[htb]
\centering
\includegraphics[width=0.45\textwidth]{aux/cl_base}
\end{figure}
\end{frame}
%%
\begin{frame}
\frametitle{Starting observation}
\begin{figure}[htb]
\centering
\includegraphics[width=1\textwidth]{aux/CL_important}
\end{figure}
\end{frame}

\begin{frame}[shrink=5]
\frametitle{Nested CES}
\begin{itemize}
\item Aggregate production function
\be
y_{t} = \left(\theta_{ht}H_{t}^{\rho}+\theta_{ct}C_{t}^{\rho}\right)^{\frac{1}{\rho}}
\ee
where $(H_{t},C_{t})$ are high school and college equivalent labor aggregates of age groups $j$:
\be
H_{t} = \left[\sum_{j}\left(\alpha_{j}H_{jt}^{\eta}\right)\right]^{\frac{1}{\eta}}
\ee
\be
C_{t} = \left[\sum_{j}\left(\beta_{j}C_{jt}^{\eta}\right)\right]^{\frac{1}{\eta}}
\ee
\item The relevant elasticities are 
\be
\sigma_{A}=(1-\eta)^{-1} \qquad \sigma_{E}=(1-\rho)^{-1}
\ee
\item The dogma of perfect competition implies that
\be
\log \dfrac{w_{jt}^{c}}{w_{jt}^{h}} = \log\dfrac{\theta_{ct}}{\theta_{ht}} +   \log\dfrac{\beta_{j}}{\alpha_{j}} + (\rho-\eta)\log \dfrac{C_{t}}{H_{t}} + (\eta-1)\log\dfrac{C_{jt}}{H_{jt}}
\ee
\end{itemize} 
\end{frame}
%%
\begin{frame}[shrink=5]
\frametitle{Interpretation}
\begin{itemize}
\item It is useful to rearrange the last expression as
\be
\log \dfrac{w_{jt}^{c}}{w_{jt}^{h}} =  \log\dfrac{\theta_{ct}}{\theta_{ht}} +   \log\dfrac{\beta_{j}}{\alpha_{j}} -\dfrac{1}{\sigma_{E}}\log \dfrac{C_{t}}{H_{t}} -\dfrac{1}{\sigma_{A}}\left[\log\dfrac{C_{jt}}{H_{jt}}-\log \dfrac{C_{t}}{H_{t}}\right]
\ee
\begin{itemize}
\item What if $\eta =1$?
\medskip
\item What if $ \log\dfrac{C_{jt}}{H_{jt}}-\log \dfrac{C_{t}}{H_{t}}$ is constant over time?
\end{itemize}
\pause
\medskip
\item How to model relative supplies?
\begin{itemize}
\item Relative supply of college grads for people aged $j$ largely determined by the college enrollment decision that these individuals made when they were 20
\be
\log\dfrac{C_{jt}}{H_{jt}}=\underbrace{\lambda_{t-j}}_{\text{Cohort Effect}}+\underbrace{\phi_{j}}_{\text{Age}}
\ee
age effect is constant over-time!
\end{itemize}
\medskip
\item So we get that 
\be
\label{CL2001_FINAL}
\log \dfrac{w_{jt}^{c}}{w_{jt}^{h}} =  \underbrace{\log\dfrac{\theta_{ct}}{\theta_{ht}}}_{\text{Year Effects}} +   \underbrace{\log\dfrac{\beta_{j}}{\alpha_{j}} -\dfrac{1}{\sigma_{A}}\phi_{j}}_{\text{Age}}-\underbrace{(\dfrac{1}{\sigma_{E}}-\dfrac{1}{\sigma_{A}})\log \dfrac{C_{t}}{H_{t}}}_{\text{Year Effects}}-\underbrace{\dfrac{1}{\sigma_{A}}\lambda_{t-j}}_{\text{Cohort}}
\ee
\end{itemize} 
\end{frame}

%%
\begin{frame}
\frametitle{Identification}
\begin{itemize} 
\item If there are cohort effects $\rightarrow$ imperfect substitution across age groups $\rightarrow$ Katz and Murphy (1992) is misspecified.
\pause
\medskip
\item But how do you estimate a Nested CES?
\medskip
\pause
\item $C_{t}$ and $H_{t}$ are model based! Discipline the model to see whether the model derived supply indexes match the data
\medskip
\be
H_{t} = \left[\sum_{j}\left(\alpha_{j}H_{jt}^{\eta}\right)\right]^{\frac{1}{\eta}}
\ee
\be
C_{t} = \left[\sum_{j}\left(\beta_{j}C_{jt}^{\eta}\right)\right]^{\frac{1}{\eta}}
\ee 
\end{itemize} 
\end{frame}
%%
\begin{frame}[shrink=5]
\frametitle{Identification}
\begin{itemize} 
\item Need ($\alpha_{j},\beta_{j},\sigma_{A})$ and I can't use equation (\ref{CL2001_FINAL})
\pause
\medskip
\item Look at relative wage across two different age groups \textit{within the same education category} across years
\be
\label{CLA1}
\log \dfrac{w_{jt}^{C}}{w_{j't}^{C}} = \log \dfrac{\beta_{j}}{\beta_{j}'} + \psi_{t}^{C} -\dfrac{1}{\sigma_{A}}\left(\log \dfrac{C_{jt}}{C_{j't}}\right)
\ee
\be
\label{CLA2}
\log \dfrac{w_{jt}^{H}}{w_{j't}^{H}} = \log \dfrac{\alpha_{j}}{\alpha_{j}'}+\psi_{t}^{H}-\dfrac{1}{\sigma_{A}}\left(\log \dfrac{H_{jt}}{H_{j't}}\right)
\ee
\item Using the equation above one can obtain (upon normalization):  ($\alpha_{j},\beta_{j},\sigma_{A})$.
\medskip
\item Steps: 
\begin{enumerate}
\item Estimate (\ref{CLA1}) and (\ref{CLA2}) $\rightarrow$ ($\alpha_{j},\beta_{j},\sigma_{A})$.
\medskip
\item Use ($\alpha_{j},\beta_{j},\sigma_{A})$ to construct $(H_{t},C_{t})$.
\medskip
\item Estimate (\ref{CL2001_FINAL}).
\end{enumerate}
\medskip
\pause
\item Gimme an overid test to run.
\end{itemize} 
\end{frame}
%%
\begin{frame}
\frametitle{Facts}
\begin{figure}[htb]
\centering
\includegraphics[width=1\textwidth]{aux/cl1}
\end{figure}
\end{frame}
%%
\begin{frame}
\frametitle{USA}
\begin{figure}[htb]
\centering
\includegraphics[width=1\textwidth]{aux/clUSA}
\end{figure}
\end{frame}
%%
\begin{frame}
\frametitle{UK}
\begin{figure}[htb]
\centering
\includegraphics[width=1\textwidth]{aux/clUK}
\end{figure}
\end{frame}
%%
\begin{frame}
\frametitle{Canada}
\begin{figure}[htb]
\centering
\includegraphics[width=1\textwidth]{aux/clCA}
\end{figure}
\end{frame}
%%
%%
\begin{frame}
\frametitle{Cohort effects weaken time trends...}
\begin{figure}[htb]
\centering
\includegraphics[width=0.8\textwidth]{aux/cl2}
\end{figure}
\end{frame}
%%
\begin{frame}
\frametitle{All the college premia driven by young guys not a general rise for all age groups!}
\begin{figure}[htb]
\centering
\includegraphics[width=0.8\textwidth]{aux/cl2}
\end{figure}
\end{frame}
%%
\begin{frame}
\frametitle{Now back-out structural parameters... Begin with CL ``first-stage"}
\pause
\begin{figure}[htb]
\centering
\includegraphics[width=0.6\textwidth]{aux/first_stage_CL}
\end{figure}
\end{frame}

%%
\begin{frame}
\frametitle{Closing the circle and estimating equation (\ref{CL2001_FINAL})}
\begin{figure}[htb]
\centering
\includegraphics[width=1\textwidth]{aux/cl3}
\end{figure}
\end{frame}
%%
\begin{frame}
\frametitle{No SBTC in Canada?}
\begin{figure}[htb]
\centering
\includegraphics[width=1\textwidth]{aux/cl3}
\end{figure}
\end{frame}
%%
\begin{frame}
\frametitle{Conclusions}
\begin{itemize}
\item Age groups imperfect substitutes $\sigma_{A}\in[4;6]$.
\medskip
\item Back to magical number $\sigma_{E}\in[1.1;1.6]$.
\medskip
\item $R^{2}$ for the US is 0.95.
\medskip
\item Relative supply explains more of growth in inequality than
previously suggested.
\medskip
\item Demand / SBTC somewhat less important than it appears.
\medskip
\pause
\item But why has the supply of college graduates stop increasing?
\end{itemize}
\end{frame}
%%
\begin{frame}
\frametitle{Discussion}
\begin{itemize}
\item What explains the slowdown of college graduation rates?
\medskip
\item CL argued that cohorts at the peak of the baby boom (1950s cohorts essentially) were crowded out of college. 
\medskip
\item Cohort crowding is present in all three countries.
\medskip
\item In a separate paper CL shows that size of the cohort is able to explain 1/4 of the intercohort slowdown in educational attainment.
\medskip
\item Other theories?
\end{itemize}
\end{frame}
%%
\begin{frame}[plain]{Autor, Levy, Murnane (2003)}
\begin{itemize}
\item Response to critiques of computer use correlations
\item Routinization hypothesis: computers substitute for human in the conduct of routine tasks.
\item Computers can also complement workers involved in carrying out problem-solving and complex communication activities (“nonroutine” tasks).
\item Use Dictionary of Occupational Titles (DOT) to figure out task content
of occupations.
\item Assess how this influences wage structure following the decrease in the price of computer capital..
\end{itemize}
\end{frame}
\begin{frame}
\frametitle{The Machine's Eye}
\begin{itemize}
    \item A task is ``routine" if it can be accomplished by machines following explicit programmed rules.
    \medskip
    \item Some tasks are analytical/interactive vs. manual. 
    \medskip
    \item Examples: Bank-Teller; Taxi Driver; Economist;
\end{itemize}
\end{frame}
\begin{frame}{Manual/Non-Manual Routine/Non-Routine}
\begin{figure}[htb]
\centering
\includegraphics[width=0.8\textwidth]{aux/ALM.pdf}
\end{figure}
\end{frame} 
\begin{frame}{Three fundamental Postulates}
\begin{enumerate}
    \item Computer capital is more substitutable for human labor in carrying out routine tasks than nonroutine tasks.
    \medskip
    \item Routine and nonroutine tasks are themselves imperfect substitutes.
    \medskip
    \item Greater intensity of routine inputs increases the marginal productivity of nonroutine inputs.
\end{enumerate}
    
\end{frame}


\begin{frame}[plain]{A task-based model}

\begin{itemize}
\item Total output in the economy is a function of computer capital ($C$), routine ($L_{R}$) and non-routine tasks ($L_{N}$):
\[
Q=\left(L_{R}+C\right)^{1-\beta}L_{N}^{\beta}
\]
\item Capital elastically provided at price $\rho$ (this is the exogenous force of the model).
\item Aggregate production follows a Cobb-Douglas production where routine tasks and capital are...
\item What is the elasticity of substitution between capital/routine and non-routine tasks? Therefore non-routine tasks are...to routine tasks.
\end{itemize}
\end{frame}
\begin{frame}[shrink=10]
\frametitle{Equilibrium}
\begin{itemize}
\item Workers inelastically supplies one unit of labor.
\item These workers are endowed with heterogeneous productivities in routine and non-routine tasks: $E_{i}=[r_{i},n_{i}]$. 
\item Endogenously decide to supply R or NR labor based on comparative advantage.
\pause
\item How is the equilibrium pinned down?
\begin{itemize}
    \item Perfect substitutability of capital and non-routine task implies that 
    \begin{equation}
    w_{R}=\rho.
    \end{equation}
    \item Let $\eta_{i}=n_{i}/r_{i}$. The marginal worker who is indifferent b/w R and NR jobs has an $\eta_{i}$ such that  
    \begin{equation}
        \eta^{*}=w_{R}/w_{N}
    \end{equation}
    \item Workers with $\eta_{i}>\eta^{*}$ supply NR labor and workers with $\eta_{i}\leq\eta^{*}$ supply R labor. 
    \pause
    \item Given the above, total supply of R and NR labor is given by  $g(\eta^{*})=\sum_{i}r_{i}\mathbf{1}\{\eta_{i}\leq\eta^{*}\}$ and  $h(\eta^{*})=\sum_{i}n_{i}\mathbf{1}\{\eta_{i}>\eta^{*}\}$.
    \pause
\medskip
    \item Productive efficiency requires
        \begin{equation}
       w_{R}=\dfrac{\partial Q}{\partial L_{R} }=(1-\beta)\theta^{-\beta}
    \end{equation}
      \begin{equation}
       w_{N}=\dfrac{\partial Q}{\partial L_{N} }=\beta\theta^{1-\beta}
    \end{equation}
    \begin{equation}
    \theta = \dfrac{C + g(\eta^{*})}{h(\eta ^{*})}
    \end{equation}
    \item Note that $\theta$ is the ratio of routine to non-routine tasks. 
    \pause
    \item Five equations in five unknowns $(C,w_{R},w_{N},\eta^{*},\theta)$!
\end{itemize}
\end{itemize}
\end{frame}
\begin{frame}[shrink=10]
\frametitle{Comparative Statics}
\begin{itemize}
    \item Exogenous force in the model: a \textit{decrease} in $\rho$ (price of computer capital).
    \item This is going to increase the relative demand for Routine-Tasks since
    \begin{equation}
        \dfrac{\partial \log \theta}{\partial \log \rho} = -\dfrac{1}{\beta}
    \end{equation}
    \item But is this increase in demand coming from a rise in computers or routine labor? \pause $\rightarrow$ Need to understand what is happening to supply of NR jobs (endogenous!)
    \begin{equation}
        \dfrac{\partial \log (\frac{w_{N}}{w_{R}})}{\partial \log \rho} = -\dfrac{1}{\beta}
    \end{equation}
        \begin{equation}
        \dfrac{\partial \log \eta^{*}}{\partial \log \rho} = \dfrac{1}{\beta}
    \end{equation}
    \pause
    \item The relative wage of non-routine jobs is becoming higher $\rightarrow$ more people now supply non-routine jobs $\rightarrow$ increase in demand for routine tasks will be met by computers, not people!
    \medskip
    \item Although routine labor input declines, an inflow of
computer capital more than compensates, yielding a net increase in the intensity of routine task input in production.
\end{itemize}
\end{frame}
\begin{frame}[plain]{Can we derive now predictions at the industry level?}
\pause
\begin{itemize}
\item Industry $j$'s output ($q_{j}$) is a function of 
routine ($r_{j}$) and non-routine ($n_{j}$) tasks:
\[
q_{j}=\left(r_{j}\right)^{1-\beta_{j}}n_{j}^{\beta_{j}}
\]
\item Product price is $p_{j}\left(q_{j}\right)\propto q_{j}^{-\nu}$
\item New key level of heterogeneity: $\beta_{j}$ $\rightarrow$  industry-specific factor share of nonroutine tasks 
\end{itemize}
\end{frame}

\begin{frame}[plain]{Industry-Level Demand}

\begin{itemize}
\item Industry $j$ takes the wages as given and then chooses quantities to maximize profits:

\begin{eqnarray*}
n_{j} & = & w_{N}^{-1/\nu}\left(\beta_{j}\left(1-\nu\right)\right)^{1/\nu}\left(\frac{w_{N}}{\rho}\cdot\frac{1-\beta_{j}}{\beta_{j}}\right)^{\left(\left(1-\beta_{j}\right)\left(1-\nu\right)\right)/\nu}\\
r_{j} & = & \rho^{-1/\nu}\left(\left(1-\beta_{j}\right)\left(1-\nu\right)\right)^{1/\nu}\left(\frac{w_{N}}{\rho}\cdot\frac{1-\beta_{j}}{\beta_{j}}\right)^{\left(\beta_{j}\left(\nu-1\right)\right)/\nu}
\end{eqnarray*}
\item Additional predictions from fall in price $\rho$ of routine-replacing capital
\begin{enumerate}
\medskip
\item For a given
price decline, the proportionate increase in demand for
routine task input is larger in routine-task-intensive ($\beta_j$
small) industries.
\medskip
\item A decline in the price of computer capital
also raises demand for nonroutine task input (R and NR are q-complements). This demand increase is proportionately larger in routine-task intensive
industries.
\end{enumerate}
\end{itemize}
\end{frame}
%%%
\begin{frame}{Measurement}
\begin{itemize}
    \item Labor’s Dictionary of Occupational Titles (DOT).
    \medskip
    \item Evaluates more than 12,000 highly detailed occupations along 44 objective and subjective
dimensions, including training times, physical demands and required worker aptitudes, temperaments, and interests.
\medskip
\item Information collapsed to 3 digits occupation codes (COC) $\rightarrow$ merged to CPS MORG.
\medskip
\item Micro-information then is aggregated at the industry level (140 cells). 
\medskip
\item How do jobs change overtime?
\begin{itemize}
    \item Extensive margin: Changes to the occupational distribution of employment, holding fix the task content at 1977.
    \medskip
    \item Intensive Margin: Changes in task content measures within occupations over the period 1977 to 1991. 
\end{itemize}
\medskip 
\end{itemize}
\end{frame}
%%%
%%%
\begin{frame}{Measurement}
\begin{itemize}
    \item DOT analyzes the tasks required to perform a given occupation.
    \medskip
    \item To understand, say, whether a job is non-routine ALM look at two variables in DOT evaluations: 
    \medskip
    \begin{itemize}
    \medskip
        \item Direction, Control, and Planning of activities (DCP)$\rightarrow$ Non-Routine Interactive.
        \medskip
        \item GEDMATH, measures quantitative reasoning requirements $\rightarrow$ Non-Routine Analytical.
        \medskip
        \item Similarly for Routine Analytical, Routine Manual and Non-Routine Manual. 
    \end{itemize}
    \medskip
    \item These numbers do not provide clear ordinal nor cardinal measures $\rightarrow$ DOT measures converted into percentile values corresponding to their rank in the 1960 distribution of task input.
    \medskip
    \item Analyze changes in these tasks from 1960.
\end{itemize}
\end{frame}
\begin{frame}[plain]

\begin{center}
\includegraphics[scale=0.3]{aux/trends_ALM.pdf}
\par\end{center}


\end{frame}
%%
\begin{frame}[plain]
\frametitle{Tasks seem to change: is it mainly a between vs. a within industry phenomenon?}

\begin{center}
\includegraphics[scale=0.4]{aux/within_ALM.pdf}
\par\end{center}
\end{frame}
\begin{frame}[plain]{Growth in industry computer use associated w/ shift from routine to
non-routine tasks}

\begin{center}
\includegraphics[scale=0.2]{aux/Autor_Levy_Murnane_tab3}
\par\end{center}
\pause
{\tiny{Limitation: measure of computer use comes from CPS and is only available for 80's and 90's. Use National
Income and Product Accounts (NIPA), which provides detailed
data on capital stocks across 42 major industries excluding
government.}}
\end{frame}

\begin{frame}[plain]{Capital less important than computers?}

\begin{center}
\includegraphics[scale=0.3]{aux/Autor_Levy_Murnane_tab4}
\par\end{center}

\end{frame}
\begin{frame}
\frametitle{Patterns hold also when looking within Occupations}
\begin{center}
\includegraphics[scale=0.45]{aux/ALM_within}
\par\end{center}

\end{frame}

\begin{frame}[shrink=20]
\frametitle{Does any of this matter?}
\begin{itemize}
    \item But how relevant are these shifts in tasks for a quantity like the college wage premium (a key quantity analyzed since KM 1992).
    \medskip
    \item Need to provide a mapping between tasks to demand of college graduates.
    \medskip
    \item ALM take a three-step process. First, estimate
    \begin{equation}
        CollegeShare_j = \alpha + \sum_{k=1}^{4}T_{j}^k\pi_{k} + e_j
    \end{equation}
    \medskip
    \item $CollegeShare_j\equiv$ college-equivalent share of
employment (in FTEs) in industry or occupation $j$ at the midpoint
of our samples.
\medskip
\item $T_{j}^k\equiv$ industry or occupation task
input in centiles during the same period.
\medskip
\item The coefficients $\pi_{k}$ provide an estimate of the demand of college graduates in the industry (or occupation) $j$ as a function of their underlying task. 
\pause
\medskip
\item ALM calls this a ``fixed-coefficient" specification: it ignores that industries might adjust their composition of college / non-college in response to changes in prices (effectively assuming that $\sigma$=0).
\medskip
\item In ALM words: ``This is an imperfect approximation: if the market price of nonroutine relative to routine tasks has risen, this calculation will understate demand shifts favoring nonroutine tasks and, by implication, college graduate employment".
    \end{itemize}
    
\end{frame}
\begin{frame}{Does any of this matter?}
\begin{itemize}
    \item This provides a way to isolate changes in demand for college-educated workers coming from changes in tasks: $\Delta T_{j,1970-1998}^{k}$
     \begin{equation}
       \widehat{CollegeShare_{j,1970-1998}} = \sum_{k=1}^{4}\Delta T_{j,1970-1998}\hat{\pi}_{k} 
    \end{equation}
    \item Two versions of this: one that uses actual changes in $\Delta T_{j}$ another one that uses predicted changes in $\Delta T_{j}$ predicted from computer adoption (Table 3 and 4). 
    \medskip
    \item Final step: benchmark predictions of demand for college grads coming from task-based model to the demand of college workers that would arise in K-M style (what KM calls $D_{t}$ see eq 18 of previous slide deck).
    \end{itemize}
    
\end{frame}
\begin{frame}[plain]{Changes in task content can explain non-trivial fraction of K-M demand
effect}

\includegraphics[scale=0.4]{aux/final_ALM0}
\end{frame}

\begin{frame}[plain]{Changes in task content can explain non-trivial fraction of K-M demand
effect}

\includegraphics[scale=0.4]{aux/final_ALM}
\end{frame}

\begin{frame}[plain]{Summary}

\begin{itemize}
\item Influential reboot of SBTC narrative ($\approx $10000 GS cites)
\item Task framework clarifies how computerization influences wages
\item Using DOT to measure tasks provides additional empirical content about mechanisms underlying demand side
\item Still very indirect evidence of technical change 

\begin{itemize}
\item Moving the goal posts: now we are explaining industry or occupation
level moments.
\item How important are occupations for overall wage inequality? (Firpo,
Fortin, Lemieux, 2009)
\end{itemize}
\item See also Acemoglu and Autor HOLE (2011) for a review of task-based approach in economics. 
\end{itemize}
\end{frame}
%%
\begin{frame}
\frametitle{Autor Katz Kearney (2008)}
\begin{itemize}
\item Actual title: ``Revising the Revisionists"
\medskip
\item Two claims of the Revisionists are under scrutiny:
\medskip
\begin{enumerate}
\item Growth of inequality was episodic rather than a secular phenomenom.
\medskip
\item Such episodic rise is largely explained by non-market forces (aka institutions) and labor force composition changes.
\end{enumerate}
 \medskip
\item Data: CPS (March + ORG). Series updated up to 2005. 
 \medskip
 \item Potential inconsistencies with ORG CPS due to the major restructuring of the earnings questions that occurred in 1994.
\end{itemize}
\end{frame}
%%
\begin{frame}
\frametitle{Where is the agreement? Where is the disagreement?}
\begin{figure}[htb]
\centering
\includegraphics[width=0.8\textwidth]{aux/akk1}
\end{figure}
\end{frame}
%%
\begin{frame}
\frametitle{Revisiting trends}
\begin{itemize}
\item 90/10 log wage differential, which is the difference in real log wages between the 90th and the 10th percentile of the wage distribution. 
\medskip
\item 90/50 - 50/10 log wage differential which measure upper tail and lower tail inequality
\medskip
\item College-HS log wage gap (between group inequality)
\medskip
\item 90/10, 90/50, 50/10 residual log wage gaps which are constructed after partialling out education, experience and gender. 
\medskip
\item Changes in residual inequality measures changes in inequality that are not explained by changes in the composition of the workforce as measured by these variables.
\end{itemize}
\end{frame}
%%
\begin{frame}
\frametitle{Episodic vs. Non Episodic Inequality}
\begin{figure}[htb]
\centering
\includegraphics[width=1\textwidth]{aux/akk3}
\end{figure}
\end{frame}
\begin{frame}
\frametitle{Min Wage plausible for lower tail inequality but what about upper tail?}
\begin{figure}[htb]
\centering
\includegraphics[width=0.85\textwidth]{aux/akk4}
\end{figure}
\end{frame}
%%
\begin{frame}
\frametitle{$\sigma=1.4$ even after controlling for minimum wage}
\begin{figure}[htb]
\centering
\includegraphics[width=1\textwidth]{aux/akk6}
\end{figure}
\end{frame}
%%
\begin{frame}
\frametitle{Extending KM up to 2005....}
\begin{figure}[htb]
\centering
\includegraphics[width=1\textwidth]{aux/AKK_NEW.pdf}
\end{figure}
\end{frame}
%%
\begin{frame}
\frametitle{Extending KM up to 2005....}
\begin{figure}[htb]
\centering
\includegraphics[width=1\textwidth]{aux/AKK_NEW1.pdf}
\end{figure}
\end{frame}
%%
%%
\begin{frame}
\frametitle{AKK Facts}
\begin{itemize}
\item Inequality grew at a lower pace in the 1990s, early 2000s than in the 1980s.
\medskip
\item Lower tail inequality largely explained by the minimum wage.
\medskip
\item But grow in inequality is not an episodic event! Upper tail inequality continued to increase in recent years while inequality in lower half stabilized somewhat. 
\medskip
\item Common theme: polarization. After a monotone surge in inequality from 1979-1987, upper tail inequality continued to grew while lower tail inequality flatted out or even was reduced. 
\end{itemize}
\end{frame}
\begin{frame}
\frametitle{But what explains polarization in inequality?}
\begin{itemize}
\item Not due to mechanical shifts in the composition of the workforce.
\medskip
\item Not due to minimum wages.
\medskip
\item Canonical supply-demand models also do not fully capture what is going on post 1992.
\pause
\medskip 
\item New concept: Tasked based production function (Autor-Levy-Murnane, 2003)
\begin{itemize}
\item Computerization has non-monotone effects on demand for skills throughout the earnings distribution.
\medskip
\item Sharp raise in the demand for abstract task (complementary).
\medskip
\item Reducing demand for clerical, mechanical routine tasks done by individuals in the middle of the earnings distribution (substitution).
\medskip
\item But no impact on manual tasks! (still need janitors)
\end{itemize}
\end{itemize}
\end{frame}
\begin{frame}
\frametitle{Census data + Dictionary of Occupational Titles}
\begin{figure}[htb]
\centering
\includegraphics[width=0.9\textwidth]{aux/akk10}
\end{figure}
\end{frame}
\begin{frame}
\frametitle{Smoking gun?}
\begin{figure}[htb]
\centering
\includegraphics[width=1\textwidth]{aux/akk11}
\end{figure}
\end{frame}
%%%%%%%%%%%%%%%%%%%%%%%%%%%
%%
\begin{frame}
\frametitle{Summarizing the saga...}
\begin{itemize}
\item (Almost) Unanimous consensus upon the following facts
\bigskip
\begin{itemize}
\item Upper and lower tail inequality exhibit different time series properties.
\bigskip
\item Relative supply shifts matters.
\bigskip
\item Minimum wage important for lower tail inequality
\bigskip
\item Occupational task content models does a good job in predicting polarization in inequality in the 1990s.
\bigskip
\item Observables far from explaining all changes to inequality.
\end{itemize}
\pause
\bigskip
\item Immense literature that generated several ``new" frontiers
\bigskip
\begin{itemize}
\item Minimum Wage.
\item Trade (China Shocks $\rightarrow $Autor, Dorn and Hanson)
\item Firms as ``institutions".
\item Impacts of Robots/AI.
\end{itemize}
\end{itemize}
\end{frame}
%%
\begin{frame}
\frametitle{Foundational Questions}
\begin{itemize}
\item US-centric literature (Dustmann, Ludsteck, Schonberg, 2009).
\bigskip
\item Largely based on survey data. Key revolution of the 2010s is to bring tax and or administrative records to this literature. 
\bigskip
\item Aggregate production function / cost functions? 
\begin{itemize}
\bigskip
\item How to incorporate realistic heterogeneity?
\item What are the limitations of Nested CES models?
\end{itemize}
\bigskip
\item Always assumed that labor markets are competitive.
\bigskip
\begin{itemize}
\item Maybe different degrees of imperfect competition over-time have role in explaining wage inequality? 
\bigskip
\item How to incorporate frictions in S-D-I models? See Daniel Haanwinckel JMP.
\end{itemize}
\end{itemize}
\end{frame}
\begin{frame}
\frametitle{New Frontiers}
\begin{figure}[htb]
\centering
\includegraphics[width=1\textwidth]{aux/autor_final.pdf}
\end{figure}
\end{frame}
\begin{frame}
\frametitle{New Frontiers}
\begin{figure}[htb]
\centering
\includegraphics[width=0.5\textwidth]{aux/acemoglu_final.pdf}
\end{figure}
\end{frame}
\begin{frame}
\frametitle{New Frontiers}
\begin{figure}[htb]
\centering
\includegraphics[width=0.5\textwidth]{aux/acemoglu_final2.pdf}
\end{figure}
\end{frame}
\end{document}

